% from intro-fr.tex
\oldpage[I]{7}
suffisamment familiers.
Nous pensons d'ailleurs que de cette façon, la lecture de ce Traité pourra être pour le débutant une bonne méthode lui permettant de se familiariser avec l'Algèbre commutative et l'Algèbre homologique, dont l'étude, lorsqu'elle ne s'accompagne pas d'applications tangibles, est jugée fastidieuse, voire déprimante, par un assez grand nombre.

\bigskip
\begin{center}*\quad*\quad*\end{center}
\bigskip

Il est hors de notre compétence de donner dans cette Introduction un aperçu historique, même sommaire, des notions et résultats exposés.
Le texte ne contiendra que des références jugées particulièrement utiles pour sa compréhension, et nous n'indiquerons l'origine que des résultats les plus importants.
Formellement du moins, les sujets traités dans notre ouvrage sont assez neufs, ce qui expliquera la rareté des références faites aux Pères de la Géométrie algébrique du \textsc{xix}\textsuperscript{e} siècle et du début du \textsc{xx}\textsuperscript{e} siècle, dont nous ne connaissons les travaux que par ouï-dire.
Il convient cependant de dire quelques mots ici sur les ouvrages qui ont le plus directement influencé les auteurs et contribué au développement du point de vue des schémas.
Il faut en tout premier lieu citer le travail fondamental (FAC) de J.-P.~Serre, qui a servi d'introduction à la Géométrie algébrique pour plus d'un jeune adepte (dont l'un des auteurs du présent Traité), rebuté par l'aridité des classiques \emph{Foundations} de A.~Weil~\cite{I-18}.
C'est là qu'il est démontré pour la première fois que la <<~topologie de Zariski~>> d'une variété algébrique <<~abstraite~>> est parfaitement appropriée pour lui appliquer certaines techniques de la Topologie algébrique et donner lieu notamment à une théorie cohomologique.
De plus, la définition d'une variété algébrique qui y est donnée est celle qui se prête le plus naturellement à l'extension de cette notion que nous développons ici\footnote{Ainsi que J.-P.~Serre nous l'a signalé, il convient de noter que l'idée de définir la structure de variété par la donnée d'un faisceau d'anneaux est due à H.~Cartan, qui a pris cette idée comme point de départ de sa théorie des espaces analytiques.
Bien entendu, tout comme en Géométrie algébrique, il importerait, en <<~Géométrie analytique~>>, de donner droit de cité aux éléments nilpotents dans les anneaux locaux des espaces analytiques.
Cette extension de la définition de H.~Cartan et J.-P.~Serre a été récemment abordée par H.~Grauert~\cite{I-5}, et il y a lieu d'espérer qu'un exposé systématique de Géométrie analytique dans ce cadre général verra bientôt le jour.
Il est d'ailleurs évident que les notions et techniques développées dans ce Traité gardent un sens en Géométrie analytique, bien qu'il faille s'attendre à des difficultés techniques plus considérables dans cette dernière théorie.
On peut prévoir que la Géométrie algébrique, par la simplicité de ses méthodes, pourra servir comme une sorte de modèle formel pour de futurs développements dans la théorie des espaces analytiques.}.
Serre avait d'ailleurs remarqué lui-même que la théorie cohomologique des variétés algébriques affines pouvait se transcrire sans difficulté en remplaçant les algèbres affines sur un corps par des anneaux commutatifs quelconques.
Les chapitres I et II de ce Traité et les deux premiers paragraphes du chapitre III peuvent donc être considérés, pour l'essentiel, comme des transpositions faciles, dans ce cadre élargi, des résultats principaux de (FAC) et d'un article ultérieur du même auteur~\cite{I-15}.
Nous avons aussi retiré grand profit du \emph{Séminaire de Géométrie algébrique} de C.~Chevalley~\cite{I-1}; en particulier, l'usage systématique des <<~ensembles constructibles~>> introduits par lui, s'est révélé fort utile en théorie des schémas (cf. chap.~IV).
Nous lui avons aussi emprunté l'étude des morphismes du point

